\section{Skalowalność algorytmów}
\label{sec:AnalizaSkalowalnosci}

Kolejnym elementem analizy była ocena skalowalności równoległych wariantów badanych algorytmów. 
W tym celu wykorzystano wartości przyspieszenia oraz efektywności uzyskane dla różnych liczb 
jednostek wykonawczych i~rozmiarów danych wejściowych.

Analizę przeprowadzono w dwóch ujęciach. 
W pierwszym przypadku oceniono zmiany przyspieszenia i~efektywności wynikające ze wzrostu liczby jednostek wykonawczych przy niezmiennym rozmiarze danych. 
W drugim natomiast oceniono, jak zwiększanie rozmiaru danych wejściowych wpływa na wyniki z zachowaniem tej samej liczby jednostek wykonawczych.

Szczegółową analizę przeprowadzono dla zbioru \texttt{random\_int}, natomiast wyniki dla zbioru \texttt{random\_float} posłużyły do oceny, 
czy podobne zależności można zaobserwować również dla danych zmiennoprzecinkowych.

\subsection*{Skalowalność względem liczby jednostek wykonawczych}

W pierwszej kolejności przeanalizowano skalowalność algorytmów przy rosnącej liczbie 
jednostek wykonawczych dla zbioru o stałym rozmiarze 1~000~000 elementów. 
Uwzględniono konfiguracje obejmujące 2, 4, 8, 16 i~32 jednostki wykonawcze.

Analiza przyspieszenia pozwala określić, jak zwiększanie liczby jednostek wykonawczych wpływało na wydajność poszczególnych algorytmów.
Wyniki zestawiono w tabeli~\ref{tab:scalability-speedup}.

\input{tables/chapter6/scalability_speedup}

Przebieg zmian przyspieszenia wraz ze wzrostem liczby jednostek wykonawczych przedstawiono również na rysunku~\ref{fig:speedup-cores-random-int}.
Linią przerywaną przedstawiono przebieg przyspieszenia idealnego, 
w~którym jego wartość zwiększa się proporcjonalnie do~liczby wykorzystanych jednostek wykonawczych.

\begin{figure}[htbp]
    \centering
    \includegraphics[width=0.95\textwidth]
    {images/chapter6/speedup/speedup_vs_cores_comparison_random_int.png}
    \caption{Przyspieszenie badanych algorytmów dla zbioru
    \texttt{random\_int} o rozmiarze 1~000~000 elementów w zależności od liczby jednostek wykonawczych. Źródło: opracowanie własne.}
    \label{fig:speedup-cores-random-int}
\end{figure}

Uzyskane wyniki wskazują, że zwiększanie liczby jednostek wykonawczych nie prowadziło do proporcjonalnego wzrostu przyspieszenia.
Dla każdego z badanych algorytmów najwyższa wartość przyspieszenia wystąpiła przy innej liczbie jednostek wykonawczych.

W przypadku Quick Sort najwyższe przyspieszenie, wynoszące 1,2078, uzyskano dla 2 jednostek wykonawczych. 
Dla 4 jednostek wartość przyspieszenia wyniosła 1,0787, natomiast przy większej liczbie jednostek spadła poniżej 1.
Dla 8, 16 i~32 jednostek osiągnęła odpowiednio 0,9777, 0,9788 i~0,9828, co oznacza,
że w tych konfiguracjach wariant równoległy nie uzyskiwał już przewagi czasowej nad wariantem sekwencyjnym.

Merge Sort osiągnął najwyższe przyspieszenie przy 4 jednostkach wykonawczych, dla których wyniosło ono 1,5912. 
Dla 2 oraz 8 jednostek uzyskano odpowiednio 1,4715 i~1,4684. 
Przy dalszym zwiększaniu liczby jednostek wykonawczych wartość przyspieszenia nie ulegała już istotnym zmianom.
Dla 16 i~32 jednostek wyniosła ona odpowiednio 1,4656 i~1,4807.
Uzyskane rezultaty wskazują, że zwiększanie liczby jednostek wykonawczych powyżej 4 nie przyczyniało się 
już do istotnej poprawy wydajności analizowanej implementacji.

Odmienny przebieg zaobserwowano dla Bucket Sort. 
Wartość przyspieszenia wzrastała od 1,1499 dla 2 jednostek wykonawczych do 1,3256 dla 4 oraz 1,6537 dla 8 jednostek.
Najwyższą wartość, wynoszącą 1,8011, uzyskano dla 16 jednostek wykonawczych.
Zwiększenie liczby jednostek wykonawczych do 32 doprowadziło do spadku wartości przyspieszenia do poziomu 1,4641.
W~przeciwieństwie do Quick Sort, dla którego przy 8, 16 i~32 jednostkach przyspieszenie spadło poniżej 1, 
Bucket Sort, Merge Sort oraz Sample Sort utrzymały przyspieszenie większe od 1 we wszystkich analizowanych konfiguracjach.

Spośród wszystkich analizowanych algorytmów najwyższą wartość przyspieszenia osiągnął Sample Sort.
Wartość ta zwiększyła się z 1,5594 dla 2 jednostek wykonawczych do 2,1325 dla 4 oraz osiągnęła 2,6021 przy 8 jednostkach.
Dalsze zwiększenie liczby jednostek spowodowało jednak spadek przyspieszenia.
Dla 16 jednostek wykonawczych wartość ta wynosiła 2,5242, natomiast dla 32 jednostek 1,8933.
Najkorzystniejszą konfiguracją Sample Sort spośród przebadanych okazało się zatem zastosowanie 8 jednostek wykonawczych.

Kolejną miarą wykorzystaną do oceny skalowalności była efektywność,
określająca przyspieszenie przypadające na jedną wykorzystywaną jednostkę wykonawczą.
Porównanie uzyskanych wartości dla badanych algorytmów zestawiono w tabeli~\ref{tab:scalability-efficiency}.

\input{tables/chapter6/scalability_efficiency}

Zmiany efektywności wraz ze zwiększaniem liczby jednostek wykonawczych przedstawiono na rysunku~\ref{fig:efficiency-cores-random-int}.
Wartość równa 1 oznacza efektywność idealną, w~której zwiększenie liczby jednostek wykonawczych 
skutkowałoby proporcjonalnym przyrostem przyspieszenia.

\begin{figure}[htbp]
    \centering
    \includegraphics[width=0.95\textwidth]
    {images/chapter6/efficiency/efficiency_vs_cores_comparison_random_int.png}
    \caption{Efektywność badanych algorytmów dla zbioru \texttt{random\_int} o rozmiarze 1~000~000 elementów 
    w zależności od liczby jednostek wykonawczych. Źródło: opracowanie własne.}
    \label{fig:efficiency-cores-random-int}
\end{figure}

Dla wszystkich badanych algorytmów wraz ze zwiększaniem liczby jednostek wykonawczych obserwowano spadek efektywności. 
Najwyższe wartości występowały dla konfiguracji wykorzystującej 2 jednostki.
Wynosiły one 0,7358 dla Merge Sort, 0,7797 dla Sample Sort, 0,6039 dla Quick Sort oraz 0,5750 dla Bucket Sort.

Przy 4 jednostkach wykonawczych najwyższą efektywność osiągnął Sample Sort, dla którego wyniosła ona 0,5331. 
Merge Sort uzyskał 0,3978, Bucket Sort 0,3314, natomiast Quick Sort 0,2697.
Przy 8 jednostkach wartości spadły odpowiednio do 0,3253 dla Sample Sort,
0,2067 dla Bucket Sort, 0,1835 dla Merge Sort oraz 0,1222 dla Quick Sort.

Dalsze zwiększanie liczby jednostek wykonawczych skutkowało wyraźnym obniżeniem efektywności. 
Przy 32 jednostkach najwyższą wartość uzyskał Sample Sort, wynoszącą 0,0592. 
W~przypadku Merge Sort, Bucket Sort oraz Quick Sort efektywność wynosiła odpowiednio 0,0463, 0,0458 oraz 0,0307.

Obserwowany spadek efektywności wraz ze zwiększaniem liczby jednostek wykonawczych wskazuje, 
że dodatkowe zasoby obliczeniowe nie przekładały się na proporcjonalny przyrost przyspieszenia. 
Największe różnice były widoczne dla konfiguracji wykorzystujących 16 i~32 jednostki,
dla których efektywność wszystkich badanych algorytmów pozostawała na niskim poziomie.

W przypadku Bucket Sort przyspieszenie rosło wraz ze zwiększaniem liczby jednostek wykonawczych
do konfiguracji obejmującej 16 jednostek, dla której osiągnęło najwyższą wartość 1,8011.
Dalsze zwiększenie liczby jednostek do 32 spowodowało spadek przyspieszenia do 1,4641.
Sample Sort uzyskał natomiast najwyższą wartość przyspieszenia spośród wszystkich badanych konfiguracji, 
wynoszącą około 2,60 przy 8 jednostkach wykonawczych.

Uzyskane wartości przyspieszenia i~efektywności świadczą o ograniczonej skalowalności badanych implementacji. 
Żaden z algorytmów nie zbliżył się do skalowania idealnego, a wraz ze wzrostem liczby jednostek wykonawczych 
różnica pomiędzy uzyskanym a teoretycznym przyspieszeniem ulegała zwiększeniu.

\subsection*{Wpływ rozmiaru danych na skalowalność}

Drugim analizowanym czynnikiem był wpływ rozmiaru danych wejściowych na uzyskiwane przyspieszenie. 
Analizę przeprowadzono dla konfiguracji obejmującej 32 jednostki wykonawcze, 
porównując wyniki dla zbiorów liczących 1~000, 10~000, 100~000 oraz 1~000~000 elementów.
Pozwoliło to ocenić, czy zwiększenie ilości przetwarzanych danych sprzyjało 
wykorzystaniu najbardziej rozbudowanej spośród badanych konfiguracji równoległych.

Zmiany wartości przyspieszenia w zależności od rozmiaru zbioru \texttt{random\_int} przedstawiono na rysunku~\ref{fig:speedup-size-random-int}.
Linią przerywaną oznaczono poziom przyspieszenia równy 1, przy którym czas wykonania wariantu równoległego i~sekwencyjnego jest taki sam.

\begin{figure}[htbp]
    \centering
    \includegraphics[width=0.95\textwidth]
    {images/chapter6/speedup/speedup_vs_data_size_random_int.png}
    \caption{Przyspieszenie badanych algorytmów w zależności od rozmiaru
    zbioru \texttt{random\_int} dla 32 jednostek wykonawczych. Źródło: opracowanie własne.}
    \label{fig:speedup-size-random-int}
\end{figure}

Uzyskane rezultaty pokazują, że wpływ zwiększania rozmiaru danych na przyspieszenie różnił się w zależności od analizowanego algorytmu.
Najbardziej zauważalną zmianę odnotowano dla Bucket Sort.
Dla zbioru zawierającego 1~000 elementów przyspieszenie wynosiło 0,9279,
natomiast dla 10~000 i~100~000 elementów wynosiło odpowiednio 1,0369 i~1,0219.
Dla największego badanego zbioru, zawierającego 1~000~000 elementów, wartość ta wzrosła do 1,4641.
Oznacza to, że dla Bucket Sort konfiguracja obejmująca 32 jednostki wykonawcze stawała się korzystniejsza wraz ze wzrostem rozmiaru danych.

Inną zależność zaobserwowano dla Merge Sort. 
Dla trzech mniejszych rozmiarów danych wartość przyspieszenia utrzymywała się poniżej 1 i~wynosiła około 0,76--0,89. 
Dla największego zbioru, zawierającego 1~000~000 elementów, przyspieszenie wzrosło do około 1,48,
co oznacza, że przy tej wielkości danych wariant równoległy uzyskał przewagę czasową nad wariantem sekwencyjnym.

W przypadku Quick Sort przyspieszenie również pozostawało poniżej 1 dla wszystkich analizowanych rozmiarów danych. Dla 1~000 elementów wynosiło 0,8568,
dla 10~000 elementów 0,8706, natomiast dla 100~000 i~1~000~000 elementów osiągnęło odpowiednio 0,8633 i~0,9828.
Mimo odnotowanego wzrostu wariant równoległy nie osiągnął przy 32 jednostkach wykonawczych przyspieszenia przekraczającego wartość 1.

Sample Sort charakteryzował się wartościami przyspieszenia zróżnicowanymi w zależności od rozmiaru danych. 
Dla 1~000 elementów wynosiło ono około 1,24, a~dla 10~000 elementów około 1,00. 
Przy 100~000 elementów wartość osiągnęła około 1,01, natomiast dla 1~000~000 elementów wzrosła do około 1,89.
W przypadku Sample Sort największy zbiór danych przyniósł zatem wyraźną poprawę przyspieszenia.

Porównanie uzyskanych rezultatów wskazuje, że zwiększanie rozmiaru danych wejściowych w zróżnicowanym 
stopniu oddziaływało na skalowalność badanych algorytmów. 
Przy zastosowaniu 32 jednostek wykonawczych największe przyspieszenie dla największego zbioru uzyskano dla Sample Sort i~wyniosło ono około 1,89.
Wysokie wartości odnotowano również dla Merge Sort oraz Bucket Sort, odpowiednio około 1,48 i~1,46.
W przypadku Quick Sort zwiększenie rozmiaru danych nie było wystarczające do uzyskania przyspieszenia przekraczającego wartość 1.

Zbliżone zależności zaobserwowano również dla zbioru \texttt{random\_float}. 
Wyniki uzyskane dla danych zmiennoprzecinkowych potwierdziły, że wpływ rozmiaru danych na skalowalność zależał od analizowanego algorytmu, 
a zwiększenie liczby przetwarzanych elementów nie gwarantowało poprawy przyspieszenia we wszystkich implementacjach.

Uzyskane wyniki wskazują na zróżnicowaną skalowalność badanych implementacji. 
Wraz ze zwiększaniem liczby jednostek wykonawczych we wszystkich algorytmach odnotowywano spadek efektywności, 
jednak tempo tego spadku oraz zmiany uzyskiwanego przyspieszenia różniły się pomiędzy poszczególnymi algorytmami.

Najbardziej korzystne wyniki przy większej liczbie jednostek wykonawczych uzyskały Bucket Sort oraz Sample Sort. 
W przypadku obu algorytmów zwiększanie liczby jednostek wykonawczych początkowo prowadziło do wzrostu przyspieszenia, 
jednak po osiągnięciu najlepszej konfiguracji dalsze zwiększanie poziomu równoległości powodowało jego spadek. 
Bucket Sort utrzymywał jednak przyspieszenie większe od 1 również przy 32 jednostkach wykonawczych, 
natomiast Sample Sort uzyskiwał wartości powyżej 1 przy 16 i~32 jednostkach.

Merge Sort osiągał najlepsze wyniki przy mniejszej liczbie jednostek wykonawczych, 
jednak również przy 16 i~32 jednostkach zachowywał przyspieszenie większe od 1. 
Odmienną charakterystykę wykazywał Quick Sort, dla którego zwiększanie liczby jednostek wykonawczych 
nie prowadziło do uzyskania przyspieszenia przekraczającego wartość 1. 
Oznacza to, że w przypadku tego algorytmu wariant równoległy nie uzyskiwał przewagi nad wykonaniem sekwencyjnym.

Analiza wpływu rozmiaru danych pokazała ponadto, że zwiększanie liczby przetwarzanych elementów 
nie zapewniało jednakowej poprawy skalowalności wszystkich algorytmów.
Największą korzyść odnotowano dla Sample Sort przy największym analizowanym zbiorze,  
natomiast dla pozostałych implementacji wpływ zwiększenia rozmiaru danych był zróżnicowany.